\documentstyle[%


righttag%


,11pt%


,amssymb%


,russian%               remove in Eng.


,izvru%                 Set izven% in Eng.


]{amsart}





% \arg ?





\newcommand{\V}{{\cal V}}


\newcommand{\tts}[1]{}





%\hoffset=-15mm%                  For Eng.


\setcounter{page}{13}


\god{2002}


\nomer{1~(476)} %                        Remove in Eng.


%\nomer{Vol.\,46, No.\,1}%               Set in Eng.


%\str{\pageref{begin}--\pageref{end}}%  Set in Eng.


\udk{517.9}





\author{\it С.М.\,Воронин, Ю.И.\,Мещерякова}


% NB:   ^^ remove in Eng.


\title{Аналитическая классификация типичных вырожденных


элементарных особых точек ростков голоморфных векторных полей


на комплексной плоскости}


\thanks{Работа частично поддержана Российским фондом


фундаментальных исследований (грант \No\,98-01-00821)


и Министерством образования (грант \No\,97-0-1.8-57).}





\date{}


%\pagestyle{myheadings}%         Set in Eng.


\pagestyle{plain} %              Remove in Eng.


\begin{document}


%\thispagestyle{plain}%          Set in Eng.


\maketitle


\label{begin}








\section*{Введение}





Пусть $\V$ --- класс ростков голоморфных векторных


полей в $(\Bbb C^2,0)$ с вырожденной элементарной особой точкой 0


(т.\,е. таких, что линейная часть ростка в нуле вырождена, но не все


ее собственные значения равны нулю).





Ростки $v$ и $\wt {v}$ из $\V$ называют {\it аналитически}


({\it формально}) {\it эквивалентными}, если существует локальная голоморфная


(формальная) замена координат $H$ в $(\Bbb C^2,0)$, переводящая один


росток в другой, $H'\cdot v=\wt v\circ H$.





Ростки $v$ и $\wt v$ из $\V$ называются


{\it орбитально аналитически} ({\it формально})


{\it эквивалентными}, если существует локальная голоморфная замена координат,


переводящая фазовый портрет одного ростка в фазовый портрет другого (если


существует формальная замена координат $H$ и формальный степенной ряд $k$ с


ненулевым свободным членом такие, что $H'\cdot v=k\cdot\wt v\circ H$).





Орбитальные аналитическая и формальная классификации ростков из


$\V$ хорошо известны (\cite{1}; \cite{2}, \S\,3, cс.\,29, 33). В


данной работе будет получена аналитическая классификация


типичных ростков из $\V$. Показано, что эта классификация имеет


два функциональных модуля и четыре числовых, что в два раза


превышает количество модулей орбитальной аналитической


классификации. Доказана также теорема о секториальной


нормализации.





\section{Формальная классификация}





Известно (\cite{1}; \cite{2}, \S\,3, с.\,29; \cite{4}), что


типичный росток из $\V$ формально орбитально эквивалентен одному


из ростков


$$


v_\lambda=x\frac{\partial}{\partial


x}+\frac{y^2}{1+\lambda y} \frac{\partial}{\partial y},


\quad\lambda\in{\Bbb C}.


$$





Пусть $\V_\lambda$ --- класс ростков, формально орбитально


эквивалентных $v_\lambda$.





\begin{thm}[{\series{m}\shape{n}\selectfont о формальной%


\ классификации \cite{6}}]


Каждый росток из $\V_\lambda$ формально эквивалентен одному из


ростков


$$


v_{\lambda,a,b}=v_\lambda (a+by),\quad a,b\in\Bbb C,\quad


a\ne 0;


$$


ростки $v_{\lambda,a,b}$ с различными индексами попарно


формально не эквивалентны.


\end{thm}





Таким образом, формальная классификация типичных ростков из $\V$ имеет три


числовых модуля, в то время как орбитальная формальная --- один.





Пусть $\V_{\lambda,a,b}$ --- класс ростков из $\V_\lambda$,


формально эквивалентных ростку $v_{\lambda,a,b}$. Формальную


замену координат $\wh H$, сопрягающую росток


$v\in\V_{\lambda,a,b}$ с его формальной нормальной формой


$v_{\lambda,a,b}$ (т.\,е. такую, что $\wh H'\cdot v=


v_{\lambda,a,b}\circ\wh H$), будем называть {\it формальной


нормализующей заменой координат} для $v$.





Из теоремы 1 следует, что каждый росток из $\V_{\lambda,a,b}$


аналитически эквивалентен некоторому ростку вида $


v=v_{\lambda,a,b}+o(|x|^3+|y|^3). $


Обозначим через $\V'_{\lambda,a,b}$ класс ростков из


$\V_{\lambda,a,b}$ такого вида.





\begin{defn}


Формальный степенной ряд $\wh


H=\sum\limits^\infty_{k=0}c_k(x) y^k$ с коэффициентами $c_k:U\to\Bbb C^2$,


голоморфными в некоторой окрестности нуля $U\subset\Bbb C$ (одной и


той же для всех $k$), будем называть {\it полуформальным}


отображением.


\end{defn}





Теорему 1 можно существенно усилить.





\begin{thm}


Формальная нормализующая замена любого ростка


$v\in\V'_{\lambda,a,b}$ является полуформальным отображением.


\end{thm}





\section{Секториальная нормализация}





Пусть $\beta\in (\pi/2,\pi)$, $a\in\Bbb C_*$, $\alpha=\arg a$.





$\Omega_1=\{(x,y)\in\Bbb C^2: |x|<\varepsilon,\; |y|<\varepsilon,\;


|\arg y-\alpha -\frac{\pi}{2}|<\beta \}$;





$\Omega_2=\{(x,y)\in\Bbb C^2: |x|<\varepsilon,\; |y|<\varepsilon,\;


|\arg y-\alpha-\frac{3\pi}{2}|<\beta \}$.





\noindent{}Пару областей $(\Omega_1,\Omega_2)$ будем называть {\it хорошим


покрытием} области $\{|x|<\varepsilon,\ 0<|y|<\varepsilon\}$;


параметры $\alpha$, $\beta$ и $\varepsilon$ будем называть


соответственно {\it направлением, раствором} и {\it радиусом}


хорошего покрытия.





\begin{defn}


Пусть $\Omega=U\times S$ --- секториальная область, причем


$U$ --- окрестность нуля в $\Bbb C$, $S\subset \Bbb C$ --- сектор с


вершиной в нуле. Полуформальное отображение $\wh H=\sum h_k(x)y^k$ с


голоморфными в $U$ коэффициентами будем называть {\it


асимптотическим} для голоморфного отображения $H:\Omega\to\Bbb C^2$ на


$\Omega$, если для любой частной суммы $H_n=\sum\limits^n_{k=0}h_k(x)y^k$


$$


|{H(x,y)-H_n(x,y)}|=o(y^n) \quad \text{при}\quad


(x,y)\in\Omega,\quad y\to 0.


$$


\end{defn}





\begin{thm}[{\series{m}\shape{n}\selectfont о секториальной нормализации}]


Для любого ростка


$v\in \V_{\lambda,a,b}$ и любого хорошего


покрытия $(\Omega_1,\Omega_2)$ с направлением $a$ с заданным


раствором и достаточно малым радиусом существует пара


биголоморфных отображений $H_j:\Omega_j\to \Bbb C^2$, $j=1,2$,


таких, что\/$:$


\begin{itemize}


\item[---] $H_j$ сопрягает на $\Omega_j$ росток $v$ с его формальной


нормальной формой $v_{\lambda,a,b};$


\item[---] некоторая формальная нормализующая замена $\wh H$ ростка $v$


является асимптотической для $H_j$ на $\Omega_j$, $j=1,2$.


\end{itemize}


\end{thm}





\section{Аналитическая классификация ростков 


класса $\V_{\lambda,a,b}$}





Пусть $\cal M_\lambda$ --- пространство всех троек $(c,\phi,\psi)$ таких, что


$c\in \Bbb C$; $\phi$ и $\psi$ голоморфны в $(\Bbb C,0)$;


$\phi (0)=\psi (0)=0$, $\phi'(0)=\exp(2\pi i\lambda)$.


Две тройки $(c,\phi,\psi)$ и $(\wt c,\wt \phi,\wt \psi)$ из


$\cal M_\lambda$ будем называть {\it эквивалентными}, если для некоторого


$C\in\Bbb C_*$


$$


\wt c=C\cdot c,\quad \wt \phi (z)\equiv C\phi (C^{-1}z), \quad\wt\psi


(z)=\psi(C^{-1}z).


$$





Пусть $M_\lambda$ --- пространство


классов эквивалентности из $\cal M_\lambda$.





\begin{thm}[{\series{m}\shape{n}\selectfont об аналитической 


классификации ростков из $\V'_{\lambda,a,b}$}]


Существует такое отображение


$$


m:\V'_{\lambda,a,b}\to M_\lambda,\quad m:v\mapsto m_v,


$$


что справедливы следующие утверждения\/$:$


\begin{itemize}


\item[$1^\circ$.] эквивалентность и эквимодальность$:$


$v\sim\wt v\Leftrightarrow m_v=m_{\Tilde v};$


\item[$2^\circ$.] реализация\/$:$ для любого $m\in M_\lambda$


существует такое $v\in \V'_{\lambda,a,b}$, что $m=m_v;$


\item[$3^\circ$.] аналитическая зависимость\/$:$


для любого аналитического семейства $v_\varepsilon$


ростков из $\V'_{\lambda,a,b}$


некоторые представители $\mu_\varepsilon$ модулей $m_{v_\varepsilon}$ также


образуют аналитическое семейство.


\end{itemize}


\end{thm}





Эта теорема является точным аналогом известной теоремы об


орбитальной аналитической классификации ростков из $\V_\lambda$


(\cite{1}; \cite{2}, \S\,3, с.\,33). Отметим, что количество модулей


в задаче об аналитической классификации увеличилось вдвое по


сравнению с задачей об орбитальной аналитической классификации,


что находится в хорошем соответствии с результатами работы


\cite{3}: орбитальная аналитическая классификация имеет два числовых


и один функциональный; аналитическая классификация --- четыре


числовых и два функциональных модуля.





\section{Преобразование монодромии и $t$-монодромия}





Пусть $v\in \V'_{\lambda,a,b}$; тогда вектор $\Vec


e=(1,0)$ является собственным вектором линейной части ростка $v$ в


нуле с собственным значением $a\ne 0$. Пусть $S_v$ --- сепаратриса


ростка $v$, соответствующая собственному вектору $e$, $\Gamma_v$


--- трансверсаль к $S_v$ в точке $P_0=S_v\bigcap \Gamma_v$. Пусть


$z$ --- локальный параметр на $\Gamma_v$, $z(P_0)=0$, тогда


можно отождествить $(\Gamma_v,P_0)$ с $(\Bbb C,0)$.





Заметим, что интегральные кривые ограничения $v|_{S_v}$ ростка


$v$ на сепаратрису $S_v$ имеют один и тот же период $T=2\pi i/a$.


Значит, любая интегральная кривая $f(t)$ поля $v$ с начальными


условиями $f(t_0)=P\in \Gamma_v$, достаточно близкими к $P_0$,


также пересекает трансверсаль $\Gamma_v$ в некоторый момент


времени $t_1$, близкий к


$$


t_0+T : f(t_1)=P_1,\quad |t_1-t_0-T|<0,1|T|.


$$





Отображение $\delta_v :(\Bbb C,0)\times \Bbb C \to (\Bbb C,0)\times \Bbb C$,


определенное равенством $\delta_v(z(P),t_0)=(z(P_1),t_1)$, будем


называть {\it отображением t-монодромии } ростка $v$,


соответствующим трансверсали $\Gamma_v$ и имеющим специальный вид


$$


\delta_v:(z,t)\mapsto (\Delta_v(z),t+A_v(z)),


$$


где $\Delta_v$


и $A_v$ голоморфны в $(\Bbb C,0)$ (отображения такого вида, в


соответствии с \cite{3}, \cite{5}, являются {\it $t$-сдвигами}).


Отметим, что первая компонента $\Delta_v$ преобразования


$t$-монодромии является обычным преобразованием монодромии


трансверсали $\Gamma_v$, соответствующим однократному обходу в


положительном направлении выколотой точки 0 на сепаратрисе.


\medskip





{\bf Пример.}


Обозначим через $\delta_{\lambda,a,b}$ отображение


$t$-монодромии нормальной формы $v_{\lambda,a,b}$. Прямые


выкладки показывают, что


$$


\delta_{\lambda,a,b}:(z,t)\mapsto


(\Delta_{\lambda,a,b}(z),t+A_{\lambda,a,b}(z)),


$$


где $\Delta_{\lambda,a,b}(z)=g^{2\pi i}_{w(z)}$ --- сдвиг за время


$2\pi i$ вдоль фазовых кривых поля $w(z)=\frac{z^2}{1+\lambda


z}\frac{\partial}{\partial z}$, а $A_{\lambda,a,b}(z)=\frac{2\pi


i}{a}-2\pi i\frac{b}{a^z}z+{\cal O}(z^2)$.


\medskip





В соответствии с \cite{3}, \cite{5} два $t$-сдвига $\delta$ и


$\wt \delta$ назовем аналитически (формально) эквивалентными,


если существует сопрягающий их обратимый голоморфный (формальный)


$t$--сдвиг $ h:h\circ \delta=\wt \delta \circ h. $





\begin{thm}


Ростки из $\V_\lambda$ аналитически эквивалентны тогда и только


тогда, когда аналитически эквивалентны их отображения


$t$-монодромии.


\end{thm}





Пусть $D_{\lambda,a,b}$ --- класс $t$-сдвигов, формально


эквивалентных $t$-сдвигу $\delta_{\lambda,a,b}$ из примера.


Аналитическая классификация $t$-сдвигов из $D_{\lambda,a,b}$


получена в \cite{3}, \cite{5}. В частности, там показано, что


аналитическая классификация $t$-сдвигов из $D_{\lambda,a,b}$


имеет четыре функциональных модуля. Но как это следует из теорем 3 и 4,


аналитическая классификация преобразований $t$-монодромии ростков


из $\V_{\lambda,a,b}$ имеет один числовой и два функциональных модуля.


Разница в количестве модулей объясняется тем, что модули в этих


двух классификационных задачах строятся разными способами, хотя


эти два типа модулей тесно связаны.





\begin{thm}


Пусть $m_v$ --- модуль ростка $v\in \V_{\lambda,a,b}$,


$\delta_v$ --- отображение $t$-монодромии ростка $v$ и $\wt


m_v$ --- модуль ростка $\delta_v$, построенный в соответствии с


$\cite{5}$. Тогда для некоторых представителей $(c,\phi,\psi)$ и


$(\phi_1,\phi_2,\psi_1,\psi_2)$ модулей $m_v$ и $\wt m_v$


справедливы представления


$$


\phi_1(z)\equiv \phi(z)-e^{2\pi


i\lambda}z,\quad \psi_1(z)\equiv \psi(z), \quad


\phi_2(z)=z+c,\quad \psi_2(z)\equiv 0.


$$


\end{thm}





{\bf Следствие.} $t$-сдвиг $\delta\in D_{\lambda,a,b}$


является преобразованием $t$-монодромии некоторого ростка


$v\in\V_{\lambda,a,b}$, если и только если для любого


представителя $(\phi_1,\phi_2,\psi_1,\psi_2)$ его модуля


$m_\delta$ выполняются условия $\psi_2\equiv 0$,


$\phi'_2(z)\equiv 1$.





\section{Доказательства}





Теоремы 1 и 2 доказываются стандартными


рассуждениями методом последовательных приближений; утверждение


единственности (теорема 1) проверяется прямыми выкладками.





Теорема 3 для векторных полей $v$ с невязкой $v-v_{\lambda,a,b}$


вида $f\cdot \frac{\partial}{\partial x}$ доказана в \cite{1} (см.


также \cite{2}, \S\,3, с.\,30). Доказательство в общем случае


аналогично. Другое доказательство теоремы 3 можно получить с


помощью операторов конечного порядка.





Модули аналитической классификации строятся по


суперпозиции секториальных нормализующих отображений так


же, как в (\cite{2}, \S\,3, с.\,30--37). Доказательство теоремы 4


получается из доказательства аналогичной теоремы о классификации


седловых резонансных ростков незначительными изменениями.





Так как


модули отображения $t$-монодромии легко выражаются через


секториальные нормализующие отображения, то утверждение теоремы 6


доказывается прямыми выкладками.





Утверждение теоремы 5 доказывается так же, как соответствующее


утверждение работы $\cite{3}$, с использованием результатов из


теоремы 6.





\begin{thebibliography}{9}





\bibitem{1}


Martinet J., Ramis J.P. {\em Probl\'eme de modules pour des


\'equations diff\'erentielles non lin\'eaires du premier ordre} //


Pabl. Math. Inst. Hautes \' Etudes Sci. -- 1982. -- V.\,55. --


P.\,63--164.





\bibitem{2}


Il'yashenko Yu. {\em Nonlinear Stokes Phenomena} // Nonlinear Stokes


Phenomena.


% Il'yashenko Yu. (Ed.).


-- Adv. in Sov. Math. --


V.\,14. -- Amer. Math. Soc., Providence, 1993. -- 287 p.





\bibitem{3}


Voronin S.M., Grinchii A.A. {\em Analytic classification of saddle resonant


singular points of holomorfic vector fields on the complex plane}


// J. of Dynamical and Сontrol Systems. -- 1996. -- V.\,2. --


\No\,1. -- P.\,1--15.





\bibitem{4}


Арнольд В.И., Ильяшенко Ю.С. {\em Обыкновенные дифференциальные


уравнения} // Итоги науки и техн. Совр. пробл. матем. --


M.: ВИНИТИ, 1985. -- T.\,1. -- С.\,7--149.





\bibitem{5}


Воронин С.М., Гринчий А.А. {\em Аналитическая классификация $t$-сдвигов}


// Челябинск. гос.\linebreak ун-т. -- Челябинск, 1996. -- 33 с. --


Деп. в ВИНИТИ 24.05.96, \No\,1689-B96. 





\bibitem{6}


Мещерякова Ю.И. {\em Формальные нормальные формы изолированных


вырожденных элементарных особых точек} // Челябинск. гос. ун-т. --


Челябинск, 1998. -- 12 c. -- Деп. в ВИНИТИ 24.09.98, \No\,2848-B98. 





\end{thebibliography}





\vskip 2cm





\post{\it Челябинский государственный}{\it Поступила}


\post{\it университет}{06.12.1999}





\label{end}


\end{document}


